% =============================================================
%  main_xrefs.tex
%
%  Macros that resolve section/figure/equation NUMBERS from the main
%  manuscript (hydrothermal_goethite_GCA.tex) for use in the SI body.
%  Requires \usepackage{xr} and \externaldocument{hydrothermal_goethite_GCA}
%  to be loaded first (see hydrothermal_goethite_GCA_SI.tex) -- xr reads
%  the main manuscript's label -> number mapping from its .aux file, so
%  the main manuscript must have been compiled at least once, and its
%  .aux must be current: recompile hydrothermal_goethite_GCA.tex, THEN
%  this SI, whenever the main text's section/figure/equation order
%  changes. If hydrothermal_goethite_GCA.aux doesn't exist yet, these
%  print "??" instead of failing outright.
%
%  Each macro is wrapped in \NoHyper...\endNoHyper (from hyperref) so it
%  prints as a plain, unlinked number -- the main manuscript and this SI
%  are submitted as separate files, so a hyperlink to a label defined in
%  the OTHER file's PDF would not resolve to anything for the reader.
%
%  IMPORTANT: \NoHyper...\endNoHyper must be wrapped in its OWN {...}
%  group here. \NoHyper turns off hyperref's linking by locally \let-ing
%  a couple of its internal commands (\hyper@linkstart among them), but
%  \endNoHyper only restores ONE of them (\hyper@link) globally -- it does
%  NOT restore the others. Without an enclosing group of our own, those
%  local \let's have no group boundary to revert at, so the first use of
%  one of these macros in the document silently disables ALL hyperlinking
%  (including ordinary in-document \ref/\eqref elsewhere in the SI) for
%  everything that follows it. Confirmed by test compile: omitting the
%  outer {...} here drops every link annotation in the whole SI PDF to
%  zero, not just these seven.
%
%  Add a new macro here (following this file's naming/pattern) rather
%  than writing \ref{}/\eqref{} directly against a main-text label in the
%  SI body, so every main-text dependency the SI has stays visible and
%  fixable in one place.
% =============================================================

% \S\ref{sec:bounding} -- "Bounding the residual helium for long steps"
% subsubsection; used with a literal "\S" already in the SI prose, so
% this macro provides only the number.
\newcommand{\MainRefBounding}{{\NoHyper\ref{sec:bounding}\endNoHyper}}

% Eq.~\eqref{eq:permlag} -- the permeation lag equation.
\newcommand{\MainEqPermlag}{{\NoHyper\eqref{eq:permlag}\endNoHyper}}

% Eq.~\eqref{eq:permode} -- the first-order ODE for n(t) with a source term.
\newcommand{\MainEqPermode}{{\NoHyper\eqref{eq:permode}\endNoHyper}}

% Eq.~\eqref{eq:steadystate} -- the steady-state permeation approximation,
% n_ss = m / (lambda*t - 1).
\newcommand{\MainEqSteadystate}{{\NoHyper\eqref{eq:steadystate}\endNoHyper}}

% Fig.~\ref{fig:arrhenius-zoom} -- zoomed Arrhenius plot, main text.
\newcommand{\MainFigArrheniusZoom}{{\NoHyper\ref{fig:arrhenius-zoom}\endNoHyper}}

% Fig.~\ref{fig:delta-zoom} -- "Delta" diffusivity plot, main text.
\newcommand{\MainFigDeltaZoom}{{\NoHyper\ref{fig:delta-zoom}\endNoHyper}}

% Fig.~\ref{fig:ampoule-permeability} -- measured vs. literature ampoule
% permeability, main text.
\newcommand{\MainFigAmpoulePermeability}{{\NoHyper\ref{fig:ampoule-permeability}\endNoHyper}}

% --- Added 2026-09-25 for the goethite dehydroxylation SI sections ---------
% \S\ref{sec:methods-goethite} -- hydrothermal goethite experiment (methods).
\newcommand{\MainRefMethodsGoethite}{{\NoHyper\ref{sec:methods-goethite}\endNoHyper}}

% \S\ref{sec:methods-sem} -- electron microscopy of thermally treated phases.
\newcommand{\MainRefMethodsSEM}{{\NoHyper\ref{sec:methods-sem}\endNoHyper}}

% \S\ref{sec:methods-kinetics} -- thermodynamic and kinetic analysis (methods).
\newcommand{\MainRefMethodsKinetics}{{\NoHyper\ref{sec:methods-kinetics}\endNoHyper}}

% \S\ref{sec:discussion-mdd} -- MDD systematics discussion.
\newcommand{\MainRefDiscussionMDD}{{\NoHyper\ref{sec:discussion-mdd}\endNoHyper}}

% \S\ref{sec:discussion-conversion} -- hydrothermal release is diffusive.
\newcommand{\MainRefDiscussionConversion}{{\NoHyper\ref{sec:discussion-conversion}\endNoHyper}}

% Fig.~\ref{fig:gt_arrhenius} -- reaction vs diffusion Arrhenius comparison.
\newcommand{\MainFigGtArrhenius}{{\NoHyper\ref{fig:gt_arrhenius}\endNoHyper}}

% Fig.~\ref{fig:gt_hold} -- 250 C hold projection test.
\newcommand{\MainFigGtHold}{{\NoHyper\ref{fig:gt_hold}\endNoHyper}}
